\section{Conclusion}\label{sec:conclusion}

Layer-adaptive KV-cache eviction is a promising axis of long-context
inference research, but its published budget rules have been heuristic,
and at least one representative formula---our own earlier one---violates
the very memory contract that justifies it, at $97.4\%$ of steps and
$+4.2\%$ mean overshoot. \sage{} rebuilds the stack around a corrected
primitive: a capped water-filling projection that provably hits the
global budget exactly at every step, with protection floors and
$O(L\log B)$ cost, wrapping decayed cumulative scoring and
dispersion-driven layer weights whose behaviour the theory explains and
the experiments isolate. Under a matched-entry protocol that removes
nominal-versus-actual confounds, on two attention architectures, with
every number regenerated from released logs, the study yields calibrated
conclusions: exact budgets are cheap; allocation should follow measured
dispersion rather than static priors, with architecture-dependent payoff;
sink and recency protection is structural, not optional; decay trades
language-modelling quality against late-context retention and
multi-instruction robustness; and history-only eviction fundamentally
cannot preserve mid-context needles at aggressive budgets---a limitation
that cache-level retention diagnostics measure directly and that
query-aware methods are positioned to address. The path forward we
consider most valuable is combining provably budget-consistent
allocation with query-aware selection under the matched-entry protocol
introduced here, at GPU scale, with the same per-step accounting released
for inspection.
